\section{Ten-positive equations: projection before full denominator recovery}
\label{sec:compact}
The two useful new local operation/degree points are \textbf{79/382} and
\textbf{80/338}, with ten independently strictly positive integer unknowns.
The whole indexed BFS943 radix67 coefficient tuple is retained, including all
four complete payloads and their finite reconstruction. Its largest port has
592,884 bits. Expanded polynomial coefficients can be larger.

\begin{center}
\begin{tabular}{llrr}
Chart & Transport weight & Binary operations & Exact degree\\\hline
$X=qw,\ Y=q^2s,\ D=q^5$ & $V^2$ &78&516\\
$X=qw,\ Y=q^2s,\ D=q^5$ & $bV$ &79&382\\
$X=q^2w,\ Y=qs,\ D=q^6$ & $V^2$ &79&456\\
$X=q^2w,\ Y=qs,\ D=q^6$ & $bV$ &80&338
\end{tabular}
\end{center}
The square-weight rows are dominated controls, not frontier claims. The
incoming \emph{Centered Residues} package already supplies $75/626,76/464$
and $77/460$; its $XY,kY$ arithmetic sharing is inherited here. The incoming
\emph{Coprime Products} package already proves the $q^2$ chart at $80/382$;
composing it with that sharing saves one operation. The substantive new chart
is $Y=qs,X=q^2w,D=q^6$, with projection before full-denominator recovery.
The older rejected chart $Y=qs,X=qw$ had only $q^3$ visibly in its first
recurrence modulus. It is not reasserted. Sign-free unshifting, the two signed
third-index exclusions, Pell doubling, and the native calibration are credited
in the copied proof sources. No historical priority or global optimum is claimed.

\paragraph{Theorem and boundary.}
At the explicitly retained ordinary Pell/native/refutation calibration,
each polynomial below has a positive zero if and only if the inherited
syntactic calculus has a finite ZFC refutation. Equivalently, ZFC proves that
absence of a positive zero is equivalent to $\operatorname{Con}(\mathrm{ZFC})$.
This neither proves consistency nor asserts that a zero exists. It is ordinary
mathematics formalizable in ZFC, not a new end-to-end Lean/Rocq derivation or
emitted object-ZFC proof. The finite verifier checks algebra and artifacts,
not the existence or absence of refutations.

\subsection{Complete polynomial definitions}
The ten inputs are
\[
 J,F,\alpha,i,T,w,s,n,\eta,\zeta\in\mathbb Z_{>0}.
\]
The fixed ports are the entire calibrated tuple $\beta,K_0,M_C,M_F$, with
$\beta\ge31$ odd, $K_0>0$, and $0<M_C<\beta<M_F<2\beta$.
Choose one complete chart:
\[
\begin{array}{c|ccccc}
 &W_0&X&Y&D&\sigma\text{ (proof only)}\\\hline
q^2&w&qw&q^2s&q^5&s/q\\
q^1&qw&q^2w&qs&q^6&s/q^2
\end{array}
\]
In both cases $X=qW_0$ and $\sigma=Y/q^3$. Define
\begin{align*}
 h&=\beta J,&q&=h+1,&C&=q-F-\alpha,\\
 k&=\eta+\zeta,&c&=kY+\eta,&A&=Y(X+1)+1,\\
 \Delta&=A^2-1,&H&=4A-5,&v&=XY^2,\\
 \ell&=(XY)(kY),&E&=1+\ell(\ell+k),\\
 e&=k-R-1,&B&=(K_0+W_0)C-F,\\
 U&=ic^2,&a&=\Delta U^2E,&F_0&=1+a,\\
 G&=EF_0,&\mathcal F&=1+4aF_0,&Q&=\Delta(\mathcal F-1),\\
 b&=cT-R,&V&=\mathcal F b-c,&K&=X(3c-2X)+2.
\end{align*}
Also set\[R=[h(q-F)+\alpha](q^2-1)+(M_C+qM_F)J.\]
Choose $L=V^2$ or $L=bV$. The final polynomial is
\begin{align}
 M&=hDH,\nonumber\\
 Z&=hD[H(n-LB)+K]-H[DB+h(e-sq^4)],\label{eq:r19carrier}\\
 P&=Q(Z^2-M^2GV^2)+(M^2G-Z^2).\label{eq:r19poly}
\end{align}
All symbols are polynomial abbreviations, not additional unknowns. Every
square and scalar multiplication is charged as a binary gate; there is no
root, division, or exponentiation instruction in the four emitted DAGs.
For proof only set $z=Z/M$, $\kappa=K/H$, $\rho=e/D$, $t=B/h$, and
$\mathcal W=1+hL$. Since $D\sigma=sq^4$ in both charts,
\begin{equation}
 z=n+\kappa+\sigma-\mathcal Wt-\rho.
 \label{eq:r19rational}
\end{equation}
None of these quotients is initially presumed integral.

\subsection{Roots, folding, and unconditional bounds}
Positive inputs imply $q\ge32$, $X,Y\ge q$, $q\mid X,Y$, $k\ge2$, $c>k$,
$A>4$, $\Delta\ge3$, $U\ge c^2$, and $E>1$. Thus $a>c$,
$\mathcal F>2c$, and $Q>c^2$. For any integral $b$, $V\ne0$: if $b=0$
then $V=-c$; otherwise $|V|\ge\mathcal F-c>1$. Identically,
\[
 \mathcal F=(1+2a)^2,\qquad Q=4\Delta^2U^2G,\qquad \gcd(E,F_0)=1.
\]
At $P=0$,
\[
 (Q-1)z^2=G(QV^2-1).
\]
For every prime dividing $G$, the two other integer factors are $-1$
modulo that prime. Rational-square valuations force $G=g^2$, $g>0$
integral. Coprimality gives $E=\tau^2$, $F_0=f_0^2$, $g=\tau f_0$.
The positive rational $y=|z|/g$ satisfies
\begin{equation}
 QV^2-(Q-1)y^2=1.\label{eq:r19adjacent}
\end{equation}
It is integral by adjacent-Pell descent: for reduced $y=r/m$, one has
$m^2\mid Q-1$, and the norm-preserving map
\[
 v'=(2Q-1)v-2(Q-1)y,\qquad y'=(2Q-1)y-2Qv
\]
from $v=|V|>1$ has nonzero integral $v'$, $|v'|<v$, and unchanged reduced
denominator of $y$. Sign reflection and descent reach $v=y=1$.
Hence $z=\varepsilon gy$ is integral, $\varepsilon\in\{1,-1\}$.

The proof integers $f=1+2a$ and $I=2Ug$ obey
\[
 f^2=\mathcal F,\quad f^2-\Delta I^2=1,\quad c^2\mid I,
 \quad Q=(\Delta I)^2.
\]
The retained signed all-index folding theorem applies to $c\le I$ and
$V\equiv-c\pmod f$, without a sign assumption on $V$. It gives
$c=\psi_A(p)$ for odd positive $p$. The first norm gives
$k=2\psi_{2XY^2+1}(n_0)$, $n_0\ge1$. At $n_0=1$, $k=2$ and
$c=2Y+1$ lies strictly between $\psi_A(1)=1$ and $\psi_A(2)=2A$.
Consequently
\begin{equation}
 n_0\ge2,\quad k\ge8XY^2+4,\quad c>8XY^3\ge8q^5,
 \quad c>X,\quad K>0.\label{eq:r19uncond}
\end{equation}
No positivity of $B,C,R,b$ has been assumed.

\subsection{Unshifted transport, including both signed exclusions}
If $b=0$, equation~\eqref{eq:r19adjacent} gives
$y^2=c^2+(c^2-1)/(Q-1)$ strictly between consecutive integers.
Thus $b\ne0$, $\operatorname{sign}V=\operatorname{sign}b$, and
$|V|>|b|$. Directly,
\[
 a^2-g^2=E[\Delta U^2E(\Delta U^2-1)-1]>0,
 \quad \mathcal F-c>3a>3g,\quad y<\tfrac32|V|.
\]
Hence $gy<V^2/2$.
For $L=bV$, suppose first that the main index is $p=3$.
Then $c=4A^2-1$. If $n_0\ge3$, then $kY>32X^2Y^5$, whereas
$A<2XY$ gives $c<16X^2Y^2$, impossible. Therefore $n_0=2$,
$k=8XY^2+4$. With $u=X+2-2Y$,
\[
 \eta=4Y^2(Xu+1)+8Y(X+1)+3-4Y.
\]
For $u\le-1$ this is negative, since $Y(X-1)>2(X+1)$; for $u\ge2$
it exceeds $k$. Thus $u=0$ or $1$, giving $2Y-X\in\{1,2\}$,
contrary to $q\mid X,Y$, $q\ge32$. Hence $p\ge5$ and
$c>c_3:=4\Delta+3$.

Put $S=\Delta I$, so $Q=S^2$. Pell classification of
\eqref{eq:r19adjacent} gives $|V|=\chi_S(j)/S$ for odd $j$;
even traces are $\pm1$ modulo $S$. The first index gives $|V|=1$.
At the third index $|V|=4Q-3$, and the two signs would imply
\[
 \mathcal F(b-4\Delta)=c-c_3\quad\text{or}\quad
 \mathcal F(b+4\Delta)=c+c_3.
\]
Both are excluded by $0<c-c_3<c+c_3<2c<\mathcal F$.
Therefore $j\ge5$, $|V|\ge16Q^2-20Q+5>8Q^2$, and
\[
 |b|>\frac{8Q^2-c}{\mathcal F}>8\mathcal F-\tfrac12
 >4\mathcal F>4g.
\]
It follows that $L=bV>0$ and $gy<3L/8<L/2$.

For either weight $e<b$, since $b-e=cT-k+1>0$, so $\rho<|b|/D$.
If $B\le-1$, then $\mathcal Wt<-L$; using $\kappa,\sigma>0$ in
\eqref{eq:r19rational} yields
\[
 n=\varepsilon gy-\kappa-\sigma+\mathcal Wt+\rho
   <L/2-L+|b|/D<0,
\]
since $D>2$. This contradicts the actual positive input $n$.
Thus $B\ge0$, hence $C>0$. Only now use the mask bound
$3q<R<q^4-2$. Equation~\eqref{eq:r19uncond} gives $c>2R$, so $b,V>0$.
The retained normalized-rank and signed auxiliary-index lemmas, with
$c^2\mid I$, $c>2R,c>2p$, and the fixed-minus residues, give $p=R$.
Mask parity yields even $q$ and odd $J$.

\subsection{Authenticate the first index before the full denominator}
Integral $z$ in~\eqref{eq:r19rational} implies
\[
 D\mid hH(e-sq^4).
\]
Since $H=4Y(X+1)-1\equiv-1\pmod q$ and $h=q-1$,
$\gcd(q,hH)=1$. As $D=q^5$ or $q^6$, this implies $q^4\mid e$.
The universal first recurrence is
\[
 k\equiv2n_0\pmod{4XY^2}.
\]
Its modulus is visibly divisible by $q^4$: $XY^2=q^5ws^2$ in the
$q^2$ chart, and $XY^2=q^4ws^2$ in the $q^1$ chart.
Therefore $q^4\mid2n_0-R-1$. Moreover
$P_0=2XY^2+1>A$, because $P_0-A=Y(2XY-X-1)>0$.
If $n_0\ge R$, monotonicity in base and index gives
$k/2\ge\psi_A(R)=c$, contrary to $c>k$. Hence
\[
 1\le n_0<R,\qquad |2n_0-R-1|<R<q^4,
\]
and necessarily $2n_0=R+1$. Returning to the full recurrence gives
\begin{equation}
 4XY^2\mid e.\label{eq:r19fullrec}
\end{equation}
In the $q^2$ chart this already gives $q^5\mid e$.
In the $q^1$ chart it does \emph{not} yet give $q^6\mid e$.
That conclusion is deliberately deferred.

\subsection{Projection first, full denominator second}
Modulo $H$, the literal numerator gives $Z\equiv hDK$. Since $H\mid M$
and $z$ is integral, $H\mid hDK$. As $D$ is a power of $q$ and
$H\equiv-1\pmod q$, cancellation is legitimate here and gives
\begin{equation}
 H\mid hK,\label{eq:r19projectionpremise}
\end{equation}
\emph{without} assuming $\rho$ integral. This is the new ordering.

Write $R=2r+1$; $r\ge48$. The retained scale-one ratio lemma gives
\[
 \xi<\frac{\psi_A(2r+1)}{\psi_{2XY^2+1}(r+1)}
 <\xi\left(1+\frac1{Y(X+1)}\right)^{2r},\qquad
 \xi=\frac{(X+1)^{2r}}{X^r}.
\]
Together with $Y<c/k<Y+1$, this implies $Y>X^r/3$ and
$X^{r+1}<3Y(X+1)<H$. The universal projection identity, odd $H$, and
\eqref{eq:r19projectionpremise} give
$H\mid h(X-2^R)(2^RX+1)$, hence also divisibility after multiplication
by $q^3$. Here is an explicit whole-product bound. Put $t_0=2^R=2\cdot4^r$.
Since $q^3h<X^4$,
\begin{align*}
 q^3h\,|(X-t_0)(t_0X+1)|
 &<X^4(X+t_0)(t_0X+1)\\
 &\le4\cdot4^rX^6+8\cdot16^rX^5<X^{r+1}.
\end{align*}
After division by $X^{r+1}$ the two terms decrease on $X\ge32,r\ge48$;
at $(32,48)$ their sum is $2^{-117}+2^{-25}<1$. Thus the entire product
is zero, and its positive factors give $X=2^R$. The universal identity then
makes $\kappa=K/H$ a positive integer. No arbitrary factor of a composite
modulus was canceled.

Since $q\mid X=2^R$ and $q\ge32$, $q=2^m$ for $m\ge5$.
The induction $2^m\ge2m$ for $m\ge2$ yields
\[
 R>3q=3\cdot2^m\ge6m,\qquad q^6\mid X.
\]
Equation~\eqref{eq:r19fullrec} now gives $q^6\mid e$ in the new chart.
Thus $\rho=e/D$ is integral in both charts. This is a property of proved
zeros, not a claim that $q^6\mid4XY^2$ for all positive inputs.

Equation~\eqref{eq:r19rational} now makes $s/q^j-\mathcal WB/h$
integral, where $j=1$ or $2$. Clear $q^jh$ and reduce first modulo
$q^j$, then modulo $h$. Since $\gcd(q,h)=1$ and $\mathcal W\equiv1\pmod h$,
\[
 q^j\mid s,\qquad h\mid B,\qquad q^3\mid Y,
 \qquad t=B/h\in\mathbb Z_{\ge0}.
\]
The missing native divisibility has been recovered, not assumed.

\subsection{Native soundness and constructive positive completeness}
Soundness now supplies positive outer data $J,F,\alpha,W_0,Y$ with
$C>0$, $X=qW_0=2^R$, $q^3\mid Y$, and
\[
 (K_0+W_0)C-F=t(q-1),\qquad t\ge0.
\]
The retained fixed-minus parity gives $R\equiv3\pmod4$; the ratio and
exact half-binomial rounding theorem identifies the native $Y$. Reconstruct
fresh source witnesses at $A_{\rm source}=Y(X+1)+2$, \emph{not} our $A$.
With $z_{\rm source}=t+1>0$, the native transport identity is exactly
\[
 (K_0+W_0)C+(q-F)-z_{\rm source}(q-1)=1.
\]
The whole indexed BFS943 calibration yields a finite syntactic ZFC refutation.
The scalar port inequalities used early are not a substitute for that decoder.

Conversely a finite refutation gives canonical data $q=2^m$, $X=2^R=qW_{\rm native}$,
$R>3q$, $R\equiv3\pmod4$, $q^3\mid Y$, and the same transport $t\ge0$.
The preceding elementary bound gives $q^6\mid X$.
Choose $w=X/q$, $s=Y/q^2$ for the $q^2$ chart; choose $w=X/q^2$, $s=Y/q$
for the $q^1$ chart. All are positive integers, and in either case
$W_0=W_{\rm native}$: the coefficient in the native transport is preserved.
Set
\[
 c=\psi_A(R),\quad k=2\psi_{P_0}((R+1)/2),\quad
 \tau=\chi_{P_0}((R+1)/2),\quad
 \eta=c-kY,\quad\zeta=k-\eta.
\]
The retained exact ratio/rounding theorem gives $0<\eta<k$.
The full recurrence and $q^6\mid X$ give integral $\rho\ge0$;
projection gives integral $0<\kappa<Xc$.

A positive multiple of the Pell matrix period modulo $c^2\tau$ supplies a
seed pair $(f_0,I_0)$ with $c^2\tau\mid I_0$. Take $i=I_0/(c^2\tau)>0$.
Then $a=\Delta I_0^2=f_0^2-1$, $F_0=f_0^2$, $G=(\tau f_0)^2$, and the
doubled pair is exactly $f=2f_0^2-1$, $I=2f_0I_0$.
The retained auxiliary-completeness lemma supplies positive $V,y$ and the
positive integer
\[
 T=\frac{V+c+Rf^2}{cf^2}.
\]
This is the actual unshifted $T$, not a possibly negative inverse shift.
For either positive weight, choose
\[
 n=gy+(1+hL)t+\rho-\kappa-\sigma.
\]
All terms are integral since canonical $\sigma=Y/q^3$ is integral.
Also $\sigma<Y<Xc$, $\kappa<Xc$, $I_0\ge c^2$, and
$f_0>(A-1)I_0>2Xc$, so $g=\tau f_0>\kappa+\sigma$.
Thus $n>0$ and its carrier is $z=gy$, proving $P=0$.
Every independent coordinate is positive. This is a fresh canonical witness
construction, not an unproved positive inverse shear.

\subsection{Exact global degrees and reproducibility}
The literal sharing is $A-1=XY+Y$ and $\ell=(XY)(kY)$.
From the incoming $q^4$ charts, remove the private $q^3$ multiplication,
change $Y$, add $D=q^4q$ or $q^4q^2$, and add $sq^4$ and $e-sq^4$.
This is net two operations in the $q^2$ chart. The $q^1$ chart additionally
splits $X$ into $W_0=qw$, $X=qW_0$, using $W_0$ in $B$.
All gates, inputs, and four fixed ports are live.

Put $q_0=\beta J$, $k_0=\eta+\zeta$, and $C_0=q_0-F-\alpha$.
The complete leading forms for the $q^2$ chart are
\begin{align*}
 [P_{V^2}]_{516}&=2^{14}q_0^{272}w^{46}s^{108}i^{20}k_0^{64}T^4C_0^2,\\
 [P_{bV}]_{382}&=-2^{10}q_0^{206}w^{34}s^{80}i^{14}k_0^{46}T^2.
\end{align*}
In the $q^1$ chart, $\deg(cT)=\deg R=4$; dropping $R$ would give an
incorrect leading form. Define
\[
 B_0=k_0sT-q_0^2(q_0-F).
\]
Then the complete forms are
\begin{align*}
 [P_{V^2}]_{456}&=2^{14}q_0^{212}w^{46}s^{104}i^{20}k_0^{60}B_0^4C_0^2,\\
 [P_{bV}]_{338}&=-2^{10}q_0^{162}w^{34}s^{78}i^{14}k_0^{44}B_0^2.
\end{align*}
Their sparse expansions have respectively $390,47,3750,274$ monomials.
The selected maximal-$T$ coefficients are $2^{14}\beta^{272}$,
$-2^{10}\beta^{206}$, $2^{14}\beta^{212}$, $-2^{10}\beta^{162}$, nonzero
at the actual $\beta>0$. These prove exact multivariate degrees, not
numerical-ray estimates. The unique top summand is $QZ^2$ for $V^2$ and
$-QM^2GV^2$ for $bV$.

The independent standard-library checker does not import the producer. It
covers every actual row with 28 free-polynomial cuts per graph, proves the
entire leading forms, checks liveness and the whole tuple, reconstructs every
coefficient byte, and rejects ten genuine algebraic corruptions per graph.
Normal and optimized Python receipts are identical. The complete unbounded
proof is also supplied in \texttt{PROOF.md}; the active local
\texttt{SOURCE\_PINS.json} authenticates all copied ordinary dependencies.
\begin{verbatim}
python3 -B build.py --check
python3 -B verify.py --output /tmp/positive-check.json
python3 -B -O verify.py --output /tmp/positive-check-O.json
python3 -B rebuild_coefficients.py
\end{verbatim}
All source intake was read-only. No remote repository was modified.
